\documentclass[12pt]{article}

\input{../../_assets/preambulo.tex}


\begin{document}

    % 1. Foto de fondo
    % 2. Título
    % 3. Encabezado Izquierdo
    % 4. Color de fondo
    % 5. Coord x del titulo
    % 6. Coord y del titulo
    % 7. Fecha

    
    \input{../../_assets/portada}
    \portadaExamen{ffccA4.jpg}{Ecuaciones\\Diferenciales II\\Examen I}{Ecuaciones Diferenciales II. Examen I}{MidnightBlue}{-8}{28}{2026}{}

    \begin{description}
        \item[Asignatura] Ecuciones Diferenciales.
        \item[Curso Académico] 2025/26.
        \item[Grado] Doble Grado en Ingeniería Informática y Matemáticas.
        \item[Grupo] Único.
        \item[Profesor] Rafael Ortega.
        \item[Descripción] Prueba 1.
        \item[Fecha] 13 de abril de 2026.
        \item[Duración] 1 hora.
    
    \end{description}
    \newpage


    % ------------------------------------
    
    \begin{ejercicio}
        Se considera el sistema:
        \begin{equation*}
            \dot{x}_1 = \sqrt{2-t}\cos x_2,\quad \dot{x}_2 = \sqrt{1+t}\cos x_1, \quad x_1(0) = 0, \quad x_2(0) = 0
        \end{equation*}
        ¿Hay unicidad de soluciones? ¿Cuál es el intervalo maximal?
    \end{ejercicio}

    \begin{ejercicio}
        Encuentra un campo continuo $X:\mathbb{R}\times\mathbb{R}\to \mathbb{R}$ y $x_0\in \mathbb{R}$ de manera que
        \begin{equation*}
            \dot{x} = X(t,x), \qquad x(0) = x_0
        \end{equation*}
        no admita solución en $\left]-2,+\infty\right[$.
    \end{ejercicio}

    \begin{ejercicio}
        ¿Hay unicidad global para el siguiente problema?
        \begin{equation*}
            \dot{x} = t^3 x^{\nicefrac{1}{3}}, \qquad x(0) = 0
        \end{equation*}
    \end{ejercicio}

    \begin{ejercicio}
        ¿Cuántas soluciones continuas $x:\mathbb{R}\to\mathbb{R}$ admite la siguiente ecuación?
        \begin{equation*}
            x(t) = \int_{0}^{1} x(s)~ds  - 2 \int_{0}^{t} x(s)~ds 
        \end{equation*}
    \end{ejercicio}

    \begin{ejercicio}
        Se considera el sistema $x'=2x(2x^2+y^2-1)$, $y'=4(2x^2+y^2-1)$.
        \begin{enumerate}
            \item[$i)$] Encuentra las soluciones constantes de este sistema.
            \item[$ii)$] Demuestra que existe una solución definida en toda la recta real con $x(0) = 0$.       
        \end{enumerate}
    \end{ejercicio}

    \newpage
    \setcounter{ejercicio}{0} % Reiniciar contador de ejercicios
    \noindent
    \textbf{Solución.}

    \begin{ejercicio}
        Se considera el sistema:
        \begin{equation*}
            \dot{x}_1 = \sqrt{2-t}\cos x_2,\quad \dot{x}_2 = \sqrt{1+t}\cos x_1, \quad x_1(0) = 0, \quad x_2(0) = 0
        \end{equation*}
        ¿Hay unicidad de soluciones? ¿Cuál es el intervalo maximal?\\

        \noindent
        Tenemos que $d=2$, con dominio $D=\left]-1,2\right[\times \mathbb{R}^2$. El campo es:
        \begin{equation*}
            X(t,x) = \left(\sqrt{2-t}\cos x_2, \sqrt{1+t}\cos x_1\right)
        \end{equation*}
        Tenemos $t_0 = 0$, $x_0 = (0,0)$ y vemos que $X\in C^1(D)$, por lo que es continuo y tenemos unicidad de soluciones. Veamos que el campo $X$ está acotado, puesto que para la norma del máximo:
        \begin{equation*}
            \|X(t,x)\|_\infty = \sup_{t\in \left]-1,2\right[} \{\sqrt{2-t}|\cos x_2|, \sqrt{1+t}|\cos x_1|\} \leq \sup_{t\in \left]-1,2\right[} \{\sqrt{2-t},\sqrt{1+t}\} \leq \sqrt{3}
        \end{equation*}
    \end{ejercicio}

    \begin{ejercicio}
        Encuentra un campo continuo $X:\mathbb{R}\times\mathbb{R}\to \mathbb{R}$ y $x_0\in \mathbb{R}$ de manera que
        \begin{equation*}
            \dot{x} = X(t,x), \qquad x(0) = x_0
        \end{equation*}
        no admita solución en $\left]-2,+\infty\right[$.\\

        \noindent
        Si tomamos el problema $\dot{x} = x^2$, $x(0) = 1$, tenemos que la solución maximal es:
        \begin{equation*}
            x(t) = \frac{1}{1-t}, \qquad t\in \left]-\infty,1\right[
        \end{equation*}
        Como la solución maximal es única, si el problema tuviera solución en $\left]-2,+\infty\right[$, entonces $x(t)$ tendría una extension continua a $t=1$, pero esto no es posible.
    \end{ejercicio}

    \begin{ejercicio}
        ¿Hay unicidad global para el siguiente problema?
        \begin{equation*}
            \dot{x} = t^3 x^{\nicefrac{1}{3}}, \qquad x(0) = 0
        \end{equation*}
        Si buscamos una solución de la forma $x(t) = ct^\alpha$, entonces:
        \begin{equation*}
            c\alpha t^{\alpha-1} = t^3 c^{\nicefrac{1}{3}} t^{\nicefrac{\alpha}{3}}
        \end{equation*}
        Tenemos entonces que:
        \begin{equation*}
            \alpha-1 = 3 + \frac{\alpha}{3}; \quad \frac{2\alpha}{3} = 4, \quad \alpha = 6, \quad c = \frac{1}{6^{\nicefrac{3}{2}}}
        \end{equation*}
        De donde:
        \begin{equation*}
            c\alpha = c^{\nicefrac{1}{3}}, \quad c6 = c^{\nicefrac{1}{3}}, \quad c^{\nicefrac{2}{3}} = \frac{1}{6}
        \end{equation*}
        Así, una solución es:
        \begin{equation*}
            x_1(t) = \left\{\begin{array}{ll}
                \frac{t^6}{6^{\nicefrac{3}{2}}} & \text{si\ } t\geq 0 \\
                 0 & \text{si\ }  t < 0
            \end{array}\right. 
        \end{equation*}
        Y otra es $x_2(t) = 0$. Por lo que no puede haber unicidad global. Otra solución es haber resuelto el problema por variables separadas.
  \end{ejercicio}

    \begin{ejercicio}
        ¿Cuántas soluciones continuas $x:\mathbb{R}\to\mathbb{R}$ admite la siguiente ecuación?
        \begin{equation*}
            x(t) = \int_{0}^{1} x(s)~ds  - 2 \int_{0}^{t} x(s)~ds 
        \end{equation*}
        Sea $x:\mathbb{R}\to \mathbb{R}$ una solución de la ecuación, el TFC nos dice que $x$ es derivable, con derivada:
        \begin{equation*}
            \dot{x}(t) = -2x(t)
        \end{equation*}
        Ahora, las soluciones de esta última son de la forma:
        \begin{equation*}
            x(t) = ce^{-2t}, \qquad t\in \mathbb{R}
        \end{equation*}
        Así, para $t=0$ tenemos a partir de la ecuación integral que:
        \begin{equation*}
            c = x(0) = \int_{0}^{1} x(s)~ds = \int_{0}^{1} ce^{-2s}~ds  = \left[-\frac{1}{2}e^{-2s}\right]_0^1 = c\left(\frac{1}{2}-\frac{1}{2e^2}\right)
        \end{equation*}
        de donde $c= 0$ y la única solución de la primera ecuación es $x(t) = 0$.
    \end{ejercicio}

    \begin{ejercicio}
        Se considera el sistema $x'=2x(2x^2+y^2-1)$, $y'=4(2x^2+y^2-1)$.
        \begin{enumerate}
            \item[$i)$] Encuentra las soluciones constantes de este sistema.
                Tenemos $D=\mathbb{R}\times \mathbb{R}^2$, buscamos soluciones de:
                \begin{align*}
                    0 &= 2x(2x^2+y^2-1) \\
                    0 &= 4(2x^2+y^2-1)
                \end{align*}
                Es equivalente al sistema $2x^2+y^2 = 1$, que es una elipse en el plano de centro el origen, radio menor $\sqrt{2}$ y radio mayor $1$.
            \item[$ii)$] Demuestra que existe una solución definida en toda la recta real con $x(0) = 0$.

                Si tomamos una solución que no esté definida en toda la recta real entonces bien explota o se acerca a la frontera, pero $\delta D =\emptyset $. Supuesto que la solución maximal tiene $\omega<\infty$, entonces la solución diverge cuando $t\to \omega$. Si esto se diera, entonces la solución tendría que cruzar la elipse, por lo que cumpliría la misma ecuación que las soluciones constantes. Como hay unicidad del problema, tendrían que ser la misma, lo que lleva a una contradicción.

                Más en detalle, si $(x(t),y(t))$ es una solución maximal definida en $\left]\alpha,\omega\right[$ con $\alpha<0<\omega$. Por reducción al absurdo, supuesto que $\omega<\infty$ el Teorema de prolongación nos dice al ser $\delta D = \emptyset $ que:
                \begin{equation*}
                    \lim_{t\to\omega}[x(t)^2 + y(t)^2]^{\nicefrac{1}{2}} = \infty
                \end{equation*}
                Si tomamos la función $f:\left]\alpha,\omega\right[\to \mathbb{R}$ dada por $f(t) = 2x(t)^2 + y(t)^2 - 1$, tenemos que:
                \begin{equation*}
                    f(0) = -1, \qquad f(t)\geq x(t)^2 + y(t)^2 - 1
                \end{equation*}
                Pero $x(t)^2 + y(t)^2 \to \infty$, por lo que $f(t)\to \infty$ si $t\to \omega$. El Teorema de Bolzano nos dice que existe $\tau \in \left]\alpha,\omega\right[$ de manera que $f(\tau) = 0$. Así, la solución $(x(t),y(t))$ cumple la condición inicial $x(\tau) = x_0, y(\tau) = y_0$ con $2x_0^2 + y_0^2 = 1$. Además, tenemos la solución constante $(x_0,y_0)$, por lo que por unicidad deberían ser las mismas, llegando a una contradicción. Así, $\omega = \infty$  y un argumento análogo serviría para $\alpha>-\infty$.
        \end{enumerate}
    \end{ejercicio}
\end{document}
